% Converted from the reference-locked manuscript.
% Final interpretation pass; operational score and inference limits follow the locked Methods.
\section*{Discussion}

TrajRisk integrates treatment-aware RF transition probabilities with a multi-scale GNN to represent hepatocyte states during APAP-induced liver injury. Four temporally defined graphs reflect the succession of metabolic and oxidative stress, injury-associated inflammation, and regenerative or recovery programs described in APAP studies.\cite{ref010,ref018,ref034,ref049} The benefit over expression-only XGBoost was localized: TrajRisk had higher average precision at 24~h, whereas XGBoost had higher global AUROC, higher AUROC at every evaluated recovery timepoint, and higher average precision at 48~h and 96~h. The 96~h advantage was over the acute-state seed, not the pseudotime-regression baseline. Reporting both aggregate and subgroup performance exposes differences that a single discrimination metric can obscure.\cite{ref050,ref051} Reduced-model comparisons support sensitivity to the temporal graphs, seed, and feature integration, but joint changes in some ablations and unequal checkpoint-selection protocols preclude assigning an independent contribution to every component. The framework's contribution is therefore a temporally structured representation and a way to examine transition-associated programs, rather than uniformly superior discrimination.

Directed transition graphs complement methods that obtain directional information from RNA velocity\cite{ref014,ref015} or clonal lineage tracing.\cite{ref016,ref052,ref053,ref054} Lineage-based inference can itself be affected by heterogeneous proliferation across a time course.\cite{ref055} TrajRisk uses an experimental sampling axis and treatment annotations, alongside learned cell-pair weights, in relation to methods that represent transition structure using velocity or learned embeddings.\cite{ref041,ref056,ref057} Two integrated datasets lacked the paired spliced and unspliced count layers required for velocity estimation.\cite{ref014,ref058} The sampling-time constraint provides a forward orientation even when later APAP-treated cells resemble untreated hepatocytes, as can occur during heterogeneous recovery.\cite{ref010} However, nearest-neighbor-derived pair labels are a modeling surrogate, not observed cell ancestry or independently validated transition probabilities. Integrating an acute-state classifier seed with expression features follows message-passing approaches used in single-cell and regulatory-network models.\cite{ref059,ref060} The present multi-scale model outperformed the development-selected single-scale comparator; that comparison does not establish superiority over every possible graph spanning the full trajectory.

Scale 3 (24--48~h) had the highest concentration of model-derived flux, while Scale 2 carried the largest total flux. Top-flux Scale 3 edges connected a small set of compensation-associated 24~h sources to predominantly lower-pseudotime 48~h targets. Persistent compensation has been described in the APAP model,\cite{ref010} and severe overdose can produce persistent p21-positive perinecrotic hepatocytes associated with impaired recovery.\cite{ref061} The concentration of flux motivates comparison with restricted transition states or bottlenecks in other dynamic systems.\cite{ref062,ref063} It supports a hypothesis-generating resolution gateway for follow-up experiments, rather than an established route followed by individual hepatocytes. Flux is the product of edge weight and source score; high-score-source enrichment is therefore partly inherent in the selection. Larger frozen-model perturbation effects at Scale 1 describe dependence of the network's predictions on that graph, not a ranking of biological necessity. Lineage tracing and interventions directed at candidate cell states or programs would be needed to test a role in recovery.

The primary human donor AUROC of 0.844 supports transfer of a mouse-trained trajectory-state score to APAP-ALF versus healthy donor discrimination in this cohort. Human nuclei contributed to unsupervised cross-species integration, but their condition labels were not used to train the seeder or GNN. No human pseudotime or trajectory was inferred. The primary self-only model transfers features, whereas the Reference-KNN sensitivity analysis also receives Gaussian-weighted messages from nearby mouse reference cells. This relates to reference mapping by transfer learning,\cite{ref064} hierarchical cell-type annotation,\cite{ref065} and cross-species alignment through shared or learned homology.\cite{ref066,ref067} Cross-species single-nucleus studies similarly distinguish shared disease programs from species-specific responses,\cite{ref068,ref069} including projection of human disease tissue onto mouse references.\cite{ref067} In our data, Reference-KNN scores and net-signature scores correlated across donors, but the signature's non-significant case--control result limits its value as a separate validation arm. Both scores share expression data, so concordance is an internal association. Using unweighted donor summaries rather than treating nuclei as biological replicates follows recommendations for avoiding pseudoreplication in single-cell comparisons.\cite{ref070,ref071,ref072}

The raw risk signature was higher in healthy donors and had AUROC 0.056 in the prespecified APAP-ALF-positive direction. Subtracting the recovery signature changed the donor ordering, but the net-signature comparison remained non-significant (p = 0.135; AUROC = 0.656) and was also non-significant in the low-nucleus-count sensitivity analysis. This result cannot be presented as a rescued or independent human validation. Mouse risk-signature genes include mature hepatocyte identity and metabolic genes such as \textit{Mat1a} and \textit{Hnf4a}. HNF4\ensuremath{\alpha}-associated function can decline during acute injury,\cite{ref073} and the MAT1A/S-adenosylmethionine axis is altered in liver disease.\cite{ref074} Injury-associated dedifferentiation\cite{ref073,ref075} and inflammatory reprogramming\cite{ref073,ref076} offer possible explanations for discordant signature direction between the mouse reference and human explants. These mechanisms were not tested here, and they do not establish why the aggregate signature reversed or exclude an imperfect cross-species mapping.

Spatial analyses test the localization of the mouse-derived gene signature, rather than spatially projected TrajRisk predictions; no cell-to-spot deconvolution was performed. Enrichment in Ben-Moshe fibrogenic injury-boundary spots is consistent with the original study's pericentral regenerative niche and extracellular-matrix induction.\cite{ref011} The risk signature was not enriched in \textit{Cyp2e1}/\textit{CYP2E1}-defined pericentral zones in either Matchett dataset. Constitutive metabolic zonation\cite{ref077,ref078,ref079} and injury-boundary annotation describe different tissue properties. Cell loss, repair-associated remodeling, and differing sampling windows\cite{ref011,ref080,ref081} could also contribute to these contrasting findings, but the present analyses do not distinguish among those explanations. Accordingly, the spatial evidence is dataset-dependent and does not establish a general pericentral risk signature. Spot-pooled p values describe differences within the profiled tissue and do not substitute for inference across independent animals or donors. Earlier spatial sampling and donor-resolved analyses could test whether the signature tracks transient injury-associated states.

Human nominal differential-expression overlap with mouse bottleneck-associated genes was 46.6\%, 46.6\%, 37.8\%, and 41.9\% across Scales 1--4. The overall 40.4\% counts distinct human symbols, rather than mapping records or an average of scale percentages. This criterion does not require matching fold-change direction between species and uses unadjusted human p values; it should therefore be interpreted as descriptive DE overlap rather than mechanistic conservation. Mouse and human responses can share some transcriptional programs while diverging elsewhere.\cite{ref082,ref083,ref084} A temporal mismatch is also plausible: experimental APAP studies resolve controlled exposure and recovery windows,\cite{ref085,ref086} whereas the human cohort comprises liver explants obtained at transplantation.\cite{ref019} The analysis cannot attribute the lower Scale 3 overlap to that mismatch, and no matched-background test establishes that any overlap exceeds chance. Donor discrimination, gene-level DE overlap, and conservation of a recovery mechanism remain separate questions.

Application beyond APAP hepatotoxicity remains to be tested. The core strategy of learning directed, condition-aware transition probabilities over discrete timepoints, decomposing a trajectory into biologically scoped temporal windows, and integrating an expression-based classifier seed with a graph neural network does not intrinsically depend on liver biology, and it is in principle applicable to other perturbation-response systems with comparable structure, namely multiple sampled timepoints, a treatment-versus-control contrast usable for initial seeding, and enough cells per transition to learn transition structure. Experimental designs of this kind are now common across longitudinal single-cell studies of perturbation, injury, and directed differentiation, including pooled genetic perturbation screens that couple defined perturbations to transcriptomic readouts,\cite{ref087} dense single-cell time courses of an acute cellular stress response,\cite{ref088} and staged multimodal atlases of directed organoid development,\cite{ref089} and harmonized perturbation-response datasets now aggregate many such experiments into common frameworks.\cite{ref090} Graph-based neural-network approaches have likewise been applied across a range of distinct single-cell problems, from cell-state modelling and imputation to gene-regulatory-network inference and prediction of unseen perturbation responses,\cite{ref024,ref091,ref092,ref093,ref094,ref095,ref096,ref097} which motivates evaluation of graph-based representations in new biological settings. Two aspects of the current implementation, however, are specific to this system and would need to be redefined elsewhere. The four temporal scales were chosen from established APAP injury and recovery biology,\cite{ref006,ref010,ref034} in which acute injury, functional compensation, recovery with fate divergence, and later resolution occur as temporally distinct phases, rather than learned automatically. The initial classifier seed also relies on a clear pair of phenotypically distinct reference states, here 0~h and 6~h, from which to initialise the seed. Automatic scale selection and seeding strategies for systems with multiple perturbations, branches, or phenotypic endpoints are natural directions for extending the framework beyond single-treatment, single-organ time courses.

\subsubsection*{Limitations of the study}

TrajRisk is trained to estimate normalized mouse pseudotime. Its within-mouse binary benchmark distinguishes cells above a pseudotime threshold, while external treatment and donor comparisons evaluate observed conditions. These analyses do not establish calibrated probabilities of subsequent injury, failure to recover, or clinical outcome; prospective outcome or lineage measurements would be needed to test such predictions.

Within-reference evaluation is transductive: training and validation cells share preconstructed graphs, and both cells' pseudotime values enter transition-pair features. The loss excludes validation cells, but target information has already influenced their graph connections. Shared expression-derived targets also limit the independence of all internal comparators. Selection of the primary checkpoint on validation AUROC further favors it over development-selected GNN ablations. Cell-level bootstrap intervals quantify resampling variability within these cohorts, rather than replication across studies or animals; within-timepoint and external pairwise improvements were not formally tested. RF pair-level cross-validation measures agreement with constructed neighbor labels, and pairs sharing cells are not independent biological transitions.

External performance varied by cohort. Reference-KNN showed only a small observed early-injury advantage and lower treatment discrimination than expression-only comparators in both later cohorts. Stronger seed and B1 correlations with \textit{Cdkn1a} similarly prevent interpreting its association as incremental benefit from TrajRisk. Because external cells contribute to joint unsupervised integration, these analyses do not test a fully frozen embedding on a new study. Prospective transfer with an independently fitted reference and additional cohorts would be needed to establish broader generalization.

Reference edges connecting external mouse cells or human nuclei to their nearest Walesky anchor cells were weighted by Gaussian similarity in scVI space rather than by the learned Random Forest transition probabilities used elsewhere in the model, because the Random Forest feature vector requires the target cell's normalized pseudotime, which is undefined for cells outside the primary trajectory. Mapping query cells onto a reference through latent-space similarity and nearest-neighbour matching is a well-established strategy in single-cell analysis.\cite{ref064,ref098,ref099} However, this means the Reference-KNN variant does not retain the learned, treatment-aware transition structure that distinguishes TrajRisk within the primary trajectory, and pseudotime imputation or joint transition inference for external cells could close this gap.

The human validation combines single-nucleus RNA-seq with mouse single-cell RNA-seq data, modalities that differ in RNA capture, transcript representation, and the relative contribution of nuclear versus cytoplasmic transcripts,\cite{ref100} which may contribute to the imperfect cross-species concordance we observed. Human APAP-ALF also differs fundamentally from the experimental mouse setting. The mouse model uses a controlled APAP dose with precisely defined post-exposure sampling intervals,\cite{ref101,ref102} whereas clinical APAP-ALF tissue is acquired at heterogeneous stages of illness and generally cannot be resolved onto the same experimentally synchronized time-since-exposure axis.\cite{ref019} Four of the ten APAP-ALF donors had fewer than 300 profiled hepatocyte nuclei, although a sensitivity analysis excluding these donors left the primary result qualitatively unchanged.

The bottleneck and flux analyses are correlational, derived from the trained model's own trajectory-state scores rather than an independent perturbational or lineage-tracing dataset, and they identify candidate transitions and model sensitivities rather than causal ones or individually fated cells. A small number of KEGG neurodegeneration terms (Parkinson, Alzheimer, and Huntington disease) enriched among Scale 1 bottleneck targets most likely reflect the extensive sharing of mitochondrial oxidative-phosphorylation and electron-transport-chain genes across these disease gene sets rather than a disease-specific association, since independent analyses have repeatedly found these neurodegenerative pathways to converge on common mitochondrial and oxidative-stress programs.\cite{ref103,ref104} The 40.4\% distinct-human-gene overlap for bottleneck-associated programs was not benchmarked against a null distribution generated from comparably expressed liver genes under the same pipeline; background-matched or permutation-based testing would provide a stronger basis for judging whether this overlap exceeds chance,\cite{ref105,ref106} and such a comparison is a natural extension of the present analysis. Finally, the binary treatment/control structure used throughout the transition model reflects the two-arm design of the datasets analysed here and would need to be generalized for studies containing multiple perturbations, doses, or conditions, an area where single-cell models that jointly represent many perturbations or doses are now emerging.\cite{ref107,ref108}


For liver injury biology, TrajRisk provides a way to organize heterogeneous hepatocyte states by sampled response phase and to prioritize compensation-associated cells and transition programs for experimental follow-up. In drug-safety research, its present role is exploratory analysis of temporally structured transcriptomic responses. Establishing prediction of adverse outcomes or treatment effects will require independent endpoints and validation beyond the current reference-based score.
